\section{Комбинаторика}
\label{sec:combinatorics}

\subsection{Основные правила подсчёта}

\textbf{Правило суммы.} Если объект можно выбрать одним из $k$ взаимно
исключающих способов, причём число вариантов равно $N_1,\ldots,N_k$, то общее
число вариантов равно
\[
    N_1+\cdots+N_k.
\]

\textbf{Правило произведения.} Если построение объекта состоит из $k$
последовательных независимых выборов с числами вариантов $N_1,\ldots,N_k$, то
число объектов равно
\[
    N_1N_2\cdots N_k.
\]

\textbf{Принцип Дирихле.} Если $N$ объектов размещены по $m$ ящикам, то хотя бы
в одном ящике находится не менее $\lceil N/m\rceil$ объектов.

\subsection{Перестановки, размещения и сочетания}

Число перестановок $n$ различных элементов:
\[
    P_n=n!.
\]
Если среди $n$ элементов имеются группы одинаковых элементов размеров
$n_1,\ldots,n_s$, где $n_1+\cdots+n_s=n$, то число различных перестановок равно
\[
    \frac{n!}{n_1!\cdots n_s!}.
\]

Число размещений без повторений из $n$ элементов по $k$:
\[
    A_n^k=\frac{n!}{(n-k)!}.
\]
Число размещений с повторениями:
\[
    \overline{A}_n^k=n^k.
\]

Число сочетаний без повторений:
\[
    C_n^k=\binom{n}{k}=\frac{n!}{k!(n-k)!}.
\]
Число сочетаний с повторениями:
\[
    \overline{C}_n^k=\binom{n+k-1}{k}.
\]
Последняя формула следует из метода ``звёзд и перегородок''.

\subsection{Биномиальные и мультиномиальные коэффициенты}

Формула бинома Ньютона:
\[
    (x+y)^n=\sum_{k=0}^{n}\binom{n}{k}x^k y^{n-k}.
\]
Основные тождества:
\[
    \binom{n}{k}=\binom{n}{n-k},
    \qquad
    \binom{n}{k}=\binom{n-1}{k-1}+\binom{n-1}{k},
\]
\[
    \sum_{k=0}^{n}\binom{n}{k}=2^n,
    \qquad
    \sum_{k=0}^{n}(-1)^k\binom{n}{k}=0\quad(n\geq1).
\]

Мультиномиальная формула:
\[
    (x_1+\cdots+x_s)^n=
    \sum_{n_1+\cdots+n_s=n}
    \frac{n!}{n_1!\cdots n_s!}
    x_1^{n_1}\cdots x_s^{n_s}.
\]

\subsection{Формула включений и исключений}

Для конечных множеств $A_1,\ldots,A_m$
\[
\left|\bigcup_{i=1}^{m}A_i\right|
=
\sum_i|A_i|
-\sum_{i<j}|A_i\cap A_j|
+\sum_{i<j<k}|A_i\cap A_j\cap A_k|
-\cdots
+(-1)^{m+1}|A_1\cap\cdots\cap A_m|.
\]
Формула применяется при подсчёте объектов, нарушающих хотя бы одно условие,
например при нахождении числа беспорядков:
\[
    D_n=n!\sum_{k=0}^{n}\frac{(-1)^k}{k!}.
\]

\subsection{Разбиения множеств}

Число разбиений $n$-элементного множества на $k$ непустых неупорядоченных
блоков называется числом Стирлинга второго рода и обозначается $S(n,k)$.
Оно удовлетворяет рекуррентному соотношению
\[
    S(n,k)=S(n-1,k-1)+kS(n-1,k).
\]
Первое слагаемое соответствует случаю, когда новый элемент образует отдельный
блок, второе --- включению его в один из $k$ существующих блоков.

Число всех разбиений множества называется числом Белла:
\[
    B_n=\sum_{k=0}^{n}S(n,k).
\]

Числа Стирлинга первого рода подсчитывают перестановки по числу циклов.
В зависимости от соглашения используют знаковую и беззнаковую версии.

\subsection{Производящие функции}

Для последовательности $a_0,a_1,\ldots$ обычной производящей функцией называется
формальный ряд
\[
    A(z)=\sum_{n=0}^{\infty}a_n z^n.
\]
Производящие функции переводят рекуррентные соотношения в алгебраические.
Например, для чисел Фибоначчи $F_0=0$, $F_1=1$,
$F_{n+2}=F_{n+1}+F_n$:
\[
    F(z)=\sum_{n\geq0}F_nz^n
    =\frac{z}{1-z-z^2}.
\]

Экспоненциальная производящая функция имеет вид
\[
    \widehat A(z)=\sum_{n=0}^{\infty}a_n\frac{z^n}{n!}
\]
и особенно удобна для подсчёта помеченных структур.

\subsection{Рекуррентные соотношения}

Линейное однородное рекуррентное соотношение с постоянными коэффициентами
\[
    a_n=c_1a_{n-1}+\cdots+c_pa_{n-p}
\]
решают с помощью характеристического многочлена
\[
    \lambda^p-c_1\lambda^{p-1}-\cdots-c_p=0.
\]
Для простых корней $\lambda_i$ решение имеет вид
\[
    a_n=\sum_i C_i\lambda_i^n.
\]
Кратному корню кратности $r$ соответствуют слагаемые
$n^j\lambda^n$, $j=0,\ldots,r-1$.

\subsection{Что важно сказать на экзамене}

Нужно различать упорядоченный и неупорядоченный выбор, наличие повторений и
ограничений. Базовый набор инструментов: правила суммы и произведения,
перестановки, размещения, сочетания, включения--исключения, числа Стирлинга и
производящие функции.
